\documentclass{article}

\usepackage{neurips_2026}

\usepackage[utf8]{inputenc}
\usepackage[T1]{fontenc}
\usepackage{hyperref}
\usepackage{url}
\usepackage{booktabs}
\usepackage{graphicx}
\usepackage{amsmath,amssymb,amsfonts,amsthm}
\usepackage{enumitem}
\usepackage{microtype}
\usepackage{xcolor}

\newtheorem{assumption}{Assumption}
\newtheorem{definition}{Definition}
\newtheorem{theorem}{Theorem}
\newtheorem{proposition}{Proposition}
\newtheorem{corollary}{Corollary}
\newtheorem{remark}{Remark}

\newcounter{algorithm}
\newenvironment{algorithm}[1][]{%
  \refstepcounter{algorithm}%
  \par\medskip
  \noindent\textbf{Algorithm~\thealgorithm}%
  \if\relax\detokenize{#1}\relax\else\ \normalfont(#1).\fi%
  \par\vspace{0.45em}%
  \normalfont\small
}{%
  \par\medskip
}

\DeclareMathOperator*{\argmin}{arg\,min}
\DeclareMathOperator*{\argmax}{arg\,max}
\DeclareMathOperator{\KL}{KL}
\DeclareMathOperator{\Occ}{Occ}
\DeclareMathOperator{\Hit}{Hit}
\DeclareMathOperator{\CVaR}{CVaR}
\DeclareMathOperator{\MI}{I}
\DeclareMathOperator{\Hent}{H}

\newcommand{\cA}{\mathcal{A}}
\newcommand{\cB}{\mathcal{B}}
\newcommand{\cD}{\mathcal{D}}
\newcommand{\cM}{\mathcal{M}}
\newcommand{\cQ}{\mathcal{Q}}
\newcommand{\cS}{\mathcal{S}}
\newcommand{\cT}{\mathcal{T}}
\newcommand{\cW}{\mathcal{W}}
\newcommand{\cX}{\mathcal{X}}
\newcommand{\cY}{\mathcal{Y}}
\newcommand{\cZ}{\mathcal{Z}}
\newcommand{\E}{\mathbb{E}}
\newcommand{\Prb}{\mathbb{P}}
\newcommand{\eps}{\varepsilon}
\newcommand{\J}{\mathcal{J}}
\newcommand{\Reg}{\mathcal{R}}
\newcommand{\reg}{\mathrm{reg}}
\newcommand{\VOI}{\mathrm{VOI}}
\newcommand{\one}{\mathbf{1}}
\newcommand{\placeholder}[2]{%
  \begin{center}
  \fbox{\begin{minipage}{#1}
  \vspace{0.4em}
  \textbf{Figure placeholder.} #2
  \vspace{0.4em}
  \end{minipage}}
  \end{center}
}

\title{How to Pick the Model:\
Minimal Diagnostic Interaction for Verifier-Backed Agentic Systems}

\author{Anonymous Authors}

\begin{document}

\maketitle

\begin{abstract}
Verifier-backed agent systems are often evaluated by whether a model can eventually solve a task.
For deployment, the sharper question is temporal and decision-theoretic: given a verifier-backed task family, a finite menu of model or harness actions, and a cheap information channel before the costly run, which action should the system choose to reach a verified success fastest, cheapest, and robustly on the current instance?
Following the time-centric view of verifiable transduction in Achille--Soatto, we treat useful task information as information that reduces search time to a verified solution rather than as an intrinsic model score.
Many such tasks admit a small pre-run interaction with the task and verifier before the long run is launched.
We formalize the router--verifier interaction as a \emph{diagnostic interface}: a standardized cheap protocol in which the verifier acts as the measuring instrument that converts limited task contact into a transcript $Z_b$ observed by a router.
The router may itself be an LLM orchestrator that controls a menu of worker models, and the chosen worker receives the expensive full run from the reset deployment state.
For any deployment loss, especially cost-aware first-hit time, we give an exact accounting identity that decomposes routed performance into the best fixed-action baseline, workload opportunity, diagnostic opportunity, and router regret.
We then define the minimum useful diagnostic interaction as the cheapest admissible measurement that recovers a target fraction of the routing value, or equivalently yields positive value after its own cost is charged.
The experimental program instantiates the framework on AutoResearch-style CIFAR-10 optimization with prompt-only Opus routing over a Haiku/Sonnet/Opus worker menu.
AutoResearch is used as a controlled member of a broader class of verifier-backed tasks: agents edit an executable artifact, receive external validation feedback, and are evaluated by time-to-verified-improvement under a fixed budget.
The main diagnostic family includes no interaction, read-only task views, verifier-only baseline probes, symmetric scratch scouts in which each candidate worker receives the same micro-budget, and adaptive diagnostic policies that decide which measurements to buy before routing.
The paper contributes a deployment-facing theory of minimal verifier-backed measurement and a benchmark protocol for testing whether diagnostic routing reduces time-to-verified-success.
\end{abstract}

\section{Introduction}
\label{sec:intro}

Suppose a builder has a verifier-backed task, a finite menu of deployable actions, and a concrete instance that must be solved under budget.
The core deployment question is not merely whether some sufficiently large model can eventually solve the task.
It is which model or harness action should be chosen \emph{now}, on this instance, before paying the full cost of a long run, in order to minimize deployment loss such as time, cost, or risk to a verified success.
This question appears in repository repair against tests, theorem proving against a checker, data-pipeline or training-loop optimization against a validation metric, and AutoResearch-style program editing where an agent modifies \texttt{train.py}, runs a bounded experiment, and keeps the edit only when validation improves~\citep{jimenez2024swebench, karpathy2026autoresearch}.
All of these are transductive and verifier-backed in the same sense: the system is given a concrete artifact-level problem, the environment provides an external success criterion, and capability is meaningful only relative to time, compute, and deployment budget.
That framing is closely aligned with the time-centric view of agents in \citet{achille2025agents}.

A scalar benchmark race is the wrong scientific object for this regime.
Two models can have similar entry success and still differ sharply in verified effort, persistence of progress, dollar cost, wall-clock cost, or robustness across workload families.
More importantly, the builder often has the option to interact \emph{a little} with the task before committing the full expensive run.
In repository repair one can run smoke tests or allow a scratch patch attempt; in theorem proving one can run a bounded proof attempt against the checker; in AutoResearch one can allow a candidate worker to make a single scratch edit and observe a short validation run.
Because the task is verifier-backed, these interactions need not be vague heuristics or hand-written feature extractors.
They can be standardized measurements that expose solver-relevant structure before the full decision is made.
This pattern is not specific to AutoResearch.
The same object appears whenever a cheap verifier interaction is available before the full run: unit-test smoke suites and scratch repairs for software, bounded checker calls for theorem proving, dry-run validation for data pipelines, short leaderboard-style submissions for ML engineering, or short training/evaluation probes for tuning tasks.
In each case the interaction is admissible only because an external verifier can turn a small amount of task contact into a router-visible signal.

The verifier is not incidental to this claim.
The same base language model can look different under different verifiers, budgets, stopping rules, and artifact interfaces.
We therefore treat an interacting model as a stochastic procedure coupled to a verifier: it proposes artifact edits, observes verifier feedback allowed by the protocol, and incurs a verifier-indexed loss.
Routing is a function of this model--verifier pair.
The verifier is both the judge of the final run and the measuring instrument that can produce task-specific information before deployment.
The diagnostic transcript does not make a model intrinsically stronger.
Rather, when it contains algorithmic information about a short path to success, it can make a non-dominant model competitive by reducing its search time on the current task.

\paragraph{Router, verifier, and worker models.}
We distinguish three roles.
The \emph{verifier} is the external measurement procedure that evaluates candidate artifacts, such as validation loss, validation accuracy, unit tests, or proof checking.
The \emph{worker models} are the models that can spend the full optimization budget on the task, for example by editing \texttt{train.py} over multiple reasoning or code-editing steps.
The \emph{router} is an LLM orchestrator: a meta-level model that observes a pre-run information state and chooses which worker model, retry policy, fallback rule, or harness action should receive the expensive run.
In the simplest AutoResearch protocol, one reasoning step is one proposed edit to \texttt{train.py} followed by a verifier run; the router does not perform the full optimization itself, but allocates that sequence of edits to a worker.
The router and workers may come from the same model family---for example Opus, Sonnet, and Haiku can all be both candidate routers and candidate workers---but their roles are distinct.
The main text can fix a single strong router first and route over a small worker menu, while later experiments can vary the router model itself.

\begin{definition}[Router--verifier diagnostic interface]
\label{def:intro-router-verifier-interface}
Fix a verifier-backed task instance $X$, verifier $V$, deployable action menu $\cA$, router model $R$, and coarse context $W$.
A router--verifier diagnostic interface is an admissible pre-deployment protocol $b$ that interacts with the task only through verifier-mediated observations, produces a transcript $\Xi_b^V(X)$, exposes a router-visible signal $Z_b=\phi_b(\Xi_b^V(X))$, and incurs diagnostic cost $c_b^V(X)$.
A router model $R$ induces a decision rule $r_{R,b}(a\mid W,Z_b)$ over actions $a\in\cA$ for the full run.
The verifier acts as a measuring instrument during $b$ and as the final judge during deployment; scratch artifacts created by $b$ are discarded unless the protocol explicitly defines them as part of the deployable action.
\end{definition}

This definition abstracts the interaction we want to measure.
The router does not need private access to the true task mode or to a hand-written recipe such as ``change the learning rate.''
It is allowed a bounded interface to the task through the verifier.
A no-edit baseline run, a smoke-test suite, a bounded proof attempt, a short validation probe, and a symmetric scratch scout are all members of the same family when they are fixed in advance, logged, costed, and separated from the final run.
The central object is therefore the deployment value of the signal produced by this interface, not the semantic label of the probe itself.

This paper asks the following question: \emph{can model or harness choice for verifier-backed agents be turned from benchmark racing into decision-theoretic performance accounting?}
Concretely, given a finite worker/action menu and a router--verifier diagnostic interface, we estimate deployment loss, diagnostic opportunity, router regret, and selection regret.
The minimum-interaction question is a subproblem of this objective: what is the cheapest admissible verifier-mediated interaction whose signal changes the action decision enough to improve net verified effort?

The paper's narrative is decision-theoretic rather than leaderboard-centric.
We do not treat routing as a mysterious capability of an especially clever LLM.
We treat it as a deployment decision problem under partial information.
The right quantities are therefore: the loss of the best fixed action, the improvement available from any observed static context, the additional value of a diagnostic interaction, and the regret of the learned router relative to the signal oracle.
Mutual information, posterior concentration, or mode-prediction accuracy are useful diagnostics, but they are secondary.
The deployment-facing question is whether a cheap interaction reduces time-to-verified-success enough to justify its cost.

\paragraph{Why AutoResearch is a good controlled family.}
AutoResearch is a particularly good experimental substrate because it isolates exactly the ingredients that matter here while staying reproducible enough for a controlled campaign.
An instance consists of an editable training program, a fixed training and evaluation budget, a verifier based on validation loss, and a nuisance seed.
The agent edits code, the environment runs the code, and success is externally certified by improved validation performance.
This is already structurally analogous to software repair against unit tests, proof search against a checker, and tuning or configuration problems where short instrumented runs expose solver-relevant difficulty before a full expensive optimization is launched.
At the same time, AutoResearch offers a disciplined controlled family: the verifier is cheap enough to probe, the artifact is concrete, counterfactual full runs can be collected for multiple models, and repeated nuisance draws give multiple deployment instances without changing the artifact class.
We use AutoResearch not as the only domain where routing should work, but as a controlled testbed for the sharper claim that verifier-backed measurement can reduce time-to-verified-improvement.
The mainline benchmark therefore uses a small family of editable starts---compact CNNs, flat MLPs, and tiny ResNets---with repeated nuisance seeds, so that every worker can be evaluated counterfactually on the same task instances.
Known family identity is treated as coarse context $W$; the main diagnostic question is whether verifier-backed pre-run contact adds value beyond that static context.

\paragraph{Contributions.}
This positioning leaves a concrete gap between algorithm selection, LLM routing, and verifier-backed agent evaluation: the field lacks an accounting language for deciding when a cheap verifier-derived signal is valuable enough to change the worker or harness action that receives the expensive run.
The paper makes four contributions.
\begin{enumerate}[leftmargin=1.35em,itemsep=0.2em]
  \item \textbf{Deployment-loss formulation.} We formulate finite-action model/harness selection under verifier-backed deployment losses, including first-hit, occupancy, and cost-aware variants.
  \item \textbf{Router--verifier diagnostic interface.} We define the verifier-mediated interaction that lets a router obtain a bounded, measurable signal about the current task before selecting the full-run action.
  \item \textbf{Routing opportunity and regret accounting.} We give an exact identity that separates best fixed performance, workload opportunity, diagnostic signal value, probe cost, and learned-router regret.
  \item \textbf{Experimental program.} We specify an AutoResearch campaign that tests metric choice, model heterogeneity, cheap-signal value, scratch-scout diagnostic sufficiency, routing recovery, minimum useful interaction budget, and selection regret, with workload-shift stress tests for robustness.
\end{enumerate}

\section{Related work and positioning}
\label{sec:related}

The closest classical ancestor is algorithm selection.
Since Rice's formulation, solver portfolios have treated performance as instance-dependent: SATzilla and AutoFolio, for example, learn to select solvers from instance features rather than report a single globally best solver~\citep{rice1976algorithm, xu2012satzilla, lindauer2015autofolio, kotthoff2016algorithm}.
That literature supplies the right high-level question, but its features are usually static descriptors computed before solving.
Our setting adds an active verifier-backed step: before selecting the full solver, the system may spend a small admissible budget to query the task and verifier, then decide whether that interaction changes the best model.

A newer line of work studies LLM routing and cascades.
FrugalGPT frames model cascades as a way to reduce cost across heterogeneous LLM APIs, while RouteLLM learns cost-quality routers from preference data~\citep{chen2023frugalgpt, ong2024routellm}.
Recent compound-system routing work extends the same idea to module allocation and multi-agent structures~\citep{chen2025optimizing, yue2025masrouter}.
These methods establish that heterogeneous model menus create economic gains, but they mostly operate on prompt-response workloads where routing signals are query features or learned response-quality predictors.
They do not ask how much task-specific verifier interaction is enough before committing an expensive agentic run.

Verifier-backed agent benchmarks provide the task substrate.
SWE-bench evaluates repository repair against tests, AutoResearch evaluates training-loop edits against validation feedback, and theorem proving or configuration tasks use checkers and validators in analogous ways~\citep{jimenez2024swebench, karpathy2026autoresearch}.
Test-time compute work asks how far a solver can be pushed with more inference, sampling, or search~\citep{snell2025testtime, brown2024monkeys, wu2025inference}.
Our question is complementary: before spending that compute, should the deployment system change which model or harness action receives the run, given a cheap verifier-derived signal?

This leaves a gap between algorithm selection, LLM routing, and verifier-backed agent evaluation.
What is missing is an accounting language that separates static context from diagnostic opportunity, probe cost, and learned-router regret.
The contributions above target this gap: the router--verifier diagnostic interface defines the admissible measurement, the action menu and verifier-indexed loss measure its deployment value, and the reusable-pilot protocol turns the idea into an experimental design.

\section{Verifier-backed deployment with diagnostic interaction}
\label{sec:setup}

\begin{definition}[Verifier-backed transductive task]
A verifier-backed transductive task is a tuple
$\cT=(\cX,\cY,V,B,\mu)$,
where $\cX$ is an instance space, $\cY$ is an artifact space, $V$ is an external verifier or verifier-derived scalar evaluation procedure, $B$ is the full deployment budget, and $\mu$ is a deployment distribution over instances.
The task is transductive because performance is measured on concrete instance-artifact pairs $(x,y)$ rather than by learning a task-independent predictor first.
\end{definition}

The verifier may be binary, as in tests or proof checking, or scalar, as in validation loss that is later thresholded into success.
The essential requirement is externality: success is not whatever the model claims, but whatever the task verifier certifies.

\begin{definition}[Worker models, LLM routers, and verifier-indexed loss]
Let $\cM=\{m_1,\ldots,m_K\}$ be a finite menu of worker models.
Given a task instance $X$, verifier $V$, and budget $B$, model $m\in\cM$ induces a stochastic interaction kernel
\[
P_m^V(dY,d\Xi,dC\mid X,B),
\]
where $Y\in\cY$ is the produced artifact, $\Xi$ is the model--verifier transcript exposed by the protocol, and $C$ is resource cost.
We call elements of $\cM$ \emph{worker models} when they are eligible to receive the full optimization run.
Let $\mathcal R$ denote a finite menu of router models.
A router model $R\in\mathcal R$ is an LLM orchestrator that observes the permitted pre-run state $(W,Z_b)$ and selects a deployable action.
A deployable action $a\in\cA$ is either a model $m$ itself or a fixed wrapper around a model, such as a retry depth, scout depth, or fallback rule.
The action incurs an integrable verifier-indexed deployment loss $\ell_V(a,X)$ derived from $(Y,\Xi,C,V)$.
In the main experimental protocol, $\cA=\cM$: the only deployment decision is which worker model receives the full verifier-backed run.
The first experimental stage may fix one strong router $R$ and evaluate its ability to choose among workers; later stages can compare router models by holding the worker menu and diagnostic interface fixed.
\end{definition}

The main theory requires only that these losses be integrable.
For readability we write $\ell(a,X)$ below, but every signal, cost, and loss is indexed by the verifier and interaction protocol.
In the experimental program we instantiate $\ell_V$ with cost-aware first-hit and occupancy losses derived from logged model--verifier trajectories.

\begin{definition}[Observed workload and admissible diagnostic interaction]
Let $W=\omega(X)\in\cW$ denote observed coarse workload information available before any probe, such as the starting program family or repository type.
Let $\cB$ be a finite family of admissible diagnostic interactions.
Each $b\in\cB$ specifies both a standardized interaction type and its diagnostic budget.
It may be nonadaptive, such as a fixed baseline verifier run, or adaptive, such as a diagnostic policy that chooses the next measurement after observing a partial transcript.
Running interaction $b$ on instance $X$ against verifier $V$ before the main run produces:
(i) a diagnostic transcript $\Xi_b^V(X)$,
(ii) a router-visible signal $Z_b=\phi_b(\Xi_b^V(X))$, and
(iii) a diagnostic cost $c_b^V(X)$.
The transcript may be verifier-only, as in a no-edit baseline evaluation, or model-coupled, as in a fixed scratch-scout panel that gives each candidate worker the same micro-budget.
Admissibility means that the diagnostic action set, budget limit, stopping rule, transcript fields, and reset rule are fixed before evaluation.
If the protocol chooses measurements adaptively, those choices must be made before the final worker is selected, must be logged and costed, and must operate on resettable scratch state whenever they modify artifacts.
\end{definition}

This definition intentionally includes both \emph{amount} and \emph{type} of interaction.
For example, $b$ may mean ``no interaction,'' ``show a read-only task view,'' ``run the baseline verifier for a small budget,'' or ``give every candidate worker one standardized scratch attempt and summarize the verifier traces.''
The protocol specifies the interaction budget and logging rules, not the semantic content of the edit, proof, patch, or experiment.
This is what makes the construction task-independent: AutoResearch scouts edit \texttt{train.py}, repository-repair scouts patch code and run tests, theorem-proving scouts propose tactics against a checker, and ML-engineering scouts make a small submission or local validation run.

\begin{definition}[Symmetric scratch scout with calibrated bid]
\label{def:scratch-scout}
A symmetric scratch scout is an admissible diagnostic interaction $b_{\mathrm{scout}}(k,B_s)$ with scout depth $k$ and scout verifier budget $B_s$.
For each worker $m\in\cM$, the protocol creates an independent scratch copy of the task, asks $m$ for a pre-scout bid about expected success and cost, allows at most $k$ model actions with verifier feedback under budget $B_s$, records the verifier deltas and resource use, and then asks for a post-scout bid.
The router-visible signal is
\[
Z_{b_{\mathrm{scout}}}
=
\left\{Z_m^{\mathrm{pre}},\Xi_{m,k}^{V,B_s},Z_m^{\mathrm{post}},C_m^{\mathrm{scout}}\right\}_{m\in\cM}.
\]
After this measurement, all scratch artifacts are discarded and the full deployment run starts from the original instance.
\end{definition}

The scratch-scout definition is the paper's preferred nontrivial measurement because it does not prescribe task-specific interventions such as changing a learning rate or toggling a scheduler.
It measures early model--task--verifier behavior under a controlled budget.
The cheaper verifier-only baseline probe remains an important ablation and a lower rung in the interaction ladder.

\begin{definition}[Adaptive diagnostic policy]
\label{def:adaptive-diagnostic}
Let $\cQ$ be a finite set of diagnostic actions, such as reading the initial artifact, running a baseline verifier probe at a specified training budget, asking for a worker bid, or running a scratch scout.
An adaptive diagnostic policy is a rule
\[
\pi_{\mathrm{diag}}(q_t\mid W,\Xi_{<t})
\]
that selects diagnostic action $q_t\in\cQ$ from the current workload context and partial transcript, stops after a bounded number of diagnostic actions or a bounded diagnostic cost, and then emits $Z_b$ to the worker router.
The cost of the adaptive diagnostic policy is included in $c_b(X)$ and therefore charged in the net objective.
\end{definition}

\paragraph{Information target.}
Let
\[
A^\star(X)\in\argmin_{a\in\cA}\E[\ell(a,X)\mid X]
\]
be the best full-run action for the instance under the chosen verifier, loss, and budget.
Before measurement, the router can only use $W$ and has uncertainty about $A^\star$ through $\Prb[A^\star\mid W]$.
After measurement, it can condition on $Z_b$.
A diagnostic interaction has task-specific decision information when
\[
\MI(A^\star;Z_b\mid W)
=
\Hent(A^\star\mid W)-\Hent(A^\star\mid W,Z_b)
>
0,
\]
but this information criterion is secondary to deployment value.
The central question is whether the measured transcript reduces cost-aware time-to-verified-success after its own cost is charged.

\paragraph{Hindsight-best worker and empirical routing capability.}
In the main worker-selection protocol, every worker model can be run counterfactually on the same held-out instance.
This lets us define the hindsight-best worker for a deployment loss $\ell$ as
\[
m^\star_\ell(X)
\in
\argmin_{m\in\cM}\ell(m,X).
\]
If router $R$ observes diagnostic signal $Z_b$ and selects worker $\hat m_R(W,Z_b)$, its instance-level routing regret is
\[
\reg_\ell(R,b;X)
=
\ell(\hat m_R(W,Z_b),X)-\min_{m\in\cM}\ell(m,X).
\]
Aggregating this regret across nuisance seeds, task modes, and repeated instances measures the empirical routing capability of the LLM orchestrator under signal $Z_b$.
An ideal router need not predict an abstract task label; it only needs to select a worker whose verified effort is close to the hindsight-best worker for the current instance and loss.

\paragraph{Connection to Achille--Soatto.}
The role of $Z_b$ is the paper's operational bridge to the Achille--Soatto view of agents as time-limited transductive solvers~\citep{achille2025agents}.
In a verifiable task, eventual solvability is too weak: an exhaustive or lucky search may eventually find a certified solution without learning anything useful about how to solve the current instance efficiently.
The relevant quantity is speed to verified success under a budget.
The diagnostic transcript $Z_b$ is therefore not treated as an intrinsic increase in model capability.
It is a finite-budget measurement channel between the current task, the candidate actions, and the verifier.
When this channel carries algorithmic information about which action has a short path to success, routing can reduce search time.
When it does not, or when the channel is too expensive, the best fixed action remains the correct deployment choice.

\paragraph{Concrete deployment losses.}
The theory below is written for a generic $\ell$, but the experimental campaign uses losses derived from relative improvement over the instance baseline.
For a scalar verifier loss $L$ where lower is better, let
$L_0(x)$ be the loss of the unedited baseline artifact and $L_t(a,x)$ the best verified loss observed by action $a$ up to full-run step $t$.
The relative improvement curve is
\[
\Delta_t(a,x)
=
\frac{L_0(x)-L_t(a,x)}{L_0(x)}.
\]
This instance-normalized criterion avoids an absolute loss target whose meaning changes across architecture, seed, verifier budget, or dataset regime.
Fix a relative-improvement threshold $\delta>0$ and horizon $H$.
Define the first threshold-hit time
\[
\tau_\delta(a,x)=\inf\{t\in\{1,\ldots,H\}:\Delta_t(a,x)\ge \delta\},
\]
with $\tau_\delta(a,x)=\infty$ if the threshold is never reached, and define threshold occupancy
\[
\Occ_{\delta,H}(a,x)=\frac{1}{H}\sum_{t=1}^H \one\{\Delta_t(a,x)\ge \delta\}.
\]
Let $C_{a,\tau_\delta\wedge H}(x)$ be the cost to first hit or horizon and $C_a^{\mathrm{full}}(x)$ the cost of the full trajectory.
Costs can be wall-clock, token/API dollars, verifier time, or a predeclared linear combination; the primary paper analysis uses wall-clock and reports token/API cost as sensitivity.
We use
\begin{align}
\ell_{\mathrm{hit}}(a,x)
&= \one\{\tau_\delta(a,x)=\infty\}+\lambda\,C_{a,\tau_\delta\wedge H}(x),
\\
\ell_{\mathrm{occ}}(a,x)
&= 1-\Occ_{\delta,H}(a,x)+\lambda\,C_a^{\mathrm{full}}(x),
\end{align}
where $\lambda$ converts cost into the same deployment scale.
The first loss measures verified entry under cost.
The second measures persistence of verified progress under cost.
For the main routing-regret analysis we also use the censored certified log-effort loss
\begin{equation}
\ell_{\log,\delta,\rho}(a,x)
=
\log C_{a,\tau_\delta\wedge H}(x)
\;+\;
\rho\,\one\{\tau_\delta(a,x)=\infty\},
\label{eq:log-effort-loss}
\end{equation}
where $\rho>0$ is the failure penalty.
This is the finite-run analogue of the certified log-effort objective used in the packed-mode theory: among successful workers it prefers lower cost-to-success; among failures it records horizon cost and adds a fixed penalty.

\section{Decision-theoretic accounting for routing}
\label{sec:accounting}

The central deployment object is the loss achieved after a diagnostic interaction and a routing decision.
We first define three oracles and then show that every learned router admits a simple exact decomposition.

\subsection{Best fixed model, workload oracle, and diagnostic oracle}

Define the best fixed-action loss
\begin{equation}
L_{\mathrm{fix}}
:=
\min_{a\in\cA} \E\big[\ell(a,X)\big].
\label{eq:lfix}
\end{equation}
This is the strongest policy available to a practitioner who chooses one interacting model or fixed harness action for the entire deployment distribution.

Define the workload oracle
\begin{equation}
L_W
:=
\E\left[\min_{a\in\cA} \E\big[\ell(a,X)\mid W\big]\right].
\label{eq:lwork}
\end{equation}
This is the Bayes-optimal policy if the builder may condition on coarse known workload identity but not on an additional probe.

For a diagnostic interaction $b\in\cB$, define the raw signal oracle
\begin{equation}
\bar L_b
:=
\E\left[\min_{a\in\cA} \E\big[\ell(a,X)\mid W,Z_b\big]\right]
\label{eq:rawlzb}
\end{equation}
and the net diagnostic oracle
\begin{equation}
L_b
:=
\E\big[c_b(X)\big]+\bar L_b.
\label{eq:netlzb}
\end{equation}
The distinction is deliberate.
$\bar L_b$ isolates the decision value of the signal itself.
$L_b$ asks whether that value survives after paying for the interaction.
Both quantities are verifier-specific: changing $V$, the budget, or the protocol changes the transcript distribution, the induced signal, and the loss surface over models.
This is exactly the point of the framework.
The paper is not asking for an intrinsic ranking of models; it asks which model is best for a concrete verifier-backed deployment.

A learned router after diagnostic interaction $b$ is a measurable policy $r_b(a\mid W,Z_b)$ over the finite action menu.
Its deployment objective is
\begin{equation}
\J_b(r_b)
:=
\E\left[c_b(X)+\sum_{a\in\cA} r_b(a\mid W,Z_b)\,\ell(a,X)\right].
\label{eq:routed-objective}
\end{equation}

\begin{definition}[Workload opportunity, diagnostic opportunity, and router regret]
For each admissible diagnostic interaction $b\in\cB$, define
\begin{align}
\Delta_W
&:= L_{\mathrm{fix}}-L_W,
\\
\bar\Delta_b
&:= L_W-\bar L_b,
\\
\Delta_b
&:= L_W-L_b = \bar\Delta_b-\E[c_b(X)],
\\
\Reg_b(r_b)
&:= \J_b(r_b)-L_b.
\end{align}
Here $\Delta_W$ is workload opportunity, $\bar\Delta_b$ is raw diagnostic opportunity, $\Delta_b$ is net diagnostic opportunity, and $\Reg_b(r_b)$ is regret relative to the diagnostic oracle.
\end{definition}

\begin{theorem}[Exact model--verifier routing identity]
\label{thm:main-identity}
For every admissible diagnostic interaction $b\in\cB$ and every learned router $r_b$,
\begin{equation}
\J_b(r_b)
=
L_{\mathrm{fix}}-\Delta_W-\Delta_b+\Reg_b(r_b).
\label{eq:main-identity}
\end{equation}
Equivalently,
\begin{equation}
\J_b(r_b)
=
L_W-\Delta_b+\Reg_b(r_b)
=
L_b+\Reg_b(r_b).
\end{equation}
\end{theorem}

\begin{proof}
The second equality is just the definition of regret.
Substituting $L_b=L_W-\Delta_b$ gives $\J_b(r_b)=L_W-\Delta_b+\Reg_b(r_b)$.
Substituting $L_W=L_{\mathrm{fix}}-\Delta_W$ yields~\eqref{eq:main-identity}.
\end{proof}

The theorem gives the paper's main accounting language.
A routed deployment can beat the best fixed action for only two reasons: coarse workload identity already matters ($\Delta_W>0$), or the diagnostic interaction adds enough verifier-backed value on top of workload identity ($\Delta_b>0$).
A learned router then keeps or loses that value through $\Reg_b(r_b)$.
This is an action-level accounting identity faithful to the operational protocol in which the router chooses an interacting model or fixed harness action rather than allocating a posterior mass directly across latent modes.

\subsection{Minimum useful diagnostic interaction}

The theorem separates two practitioner questions.
One is informational: how much routing value is in principle exposed by the signal?
The other is economic: after paying for the interaction, was the probe worth it?

\begin{definition}[Minimum sufficient and economically optimal interaction]
Fix a target recovery fraction $\eta\in(0,1]$.
The minimum sufficient diagnostic interaction is any solution of
\begin{equation}
 b^{\mathrm{suf}}_\eta
 \in
 \argmin_{b\in\cB}
 \left\{\E[c_b(X)] : \bar\Delta_b \ge \eta\,\max_{b'\in\cB}\bar\Delta_{b'}\right\}.
 \label{eq:min-sufficient}
\end{equation}
The economically optimal diagnostic interaction is any solution of
\begin{equation}
 b^{\mathrm{net}}_\star
 \in
 \argmax_{b\in\cB} \Delta_b.
 \label{eq:net-optimal}
\end{equation}
\end{definition}

The first quantity captures minimal interaction in the pure decision sense: what is the cheapest admissible probe that recovers most of the routing information that is available at all?
The second captures the deployment sense: which probe actually maximizes net value after its own cost is counted?
They need not coincide.
A richer probe can reveal more but still be uneconomical.
When the target is information rather than value, we also report the first measurably informative interaction,
\[
b^{\mathrm{info}}_{\min}
\in
\argmin_{b\in\cB}
\left\{\E[c_b(X)] : \MI(A^\star;Z_b\mid W)>0\right\},
\]
with the inequality interpreted empirically through held-out log-loss or bootstrap intervals.
This information target is useful for diagnosis, but it is not the main deployment criterion: an interaction can reduce uncertainty about the best model and still be too expensive to run.

\begin{proposition}[More informative probes help before cost is counted]
\label{prop:monotone}
Suppose $b_1,b_2\in\cB$ satisfy that $Z_{b_1}$ is measurable with respect to $Z_{b_2}$, so $b_2$ is at least as informative as $b_1$.
Then
\begin{equation}
\bar L_{b_2}\le \bar L_{b_1}
\qquad\text{and}\qquad
\bar\Delta_{b_2}\ge \bar\Delta_{b_1}.
\end{equation}
Thus richer pre-run interaction cannot hurt the signal oracle before probe cost is counted.
\end{proposition}

\begin{proof}
Conditioning on a finer signal can only enlarge the class of measurable decision rules available to the oracle.
Therefore the conditional Bayes risk under $Z_{b_2}$ is no larger than the conditional Bayes risk under $Z_{b_1}$ pointwise, and averaging proves the claim.
\end{proof}

Proposition~\ref{prop:monotone} is intentionally about raw diagnostic opportunity rather than net value.
Once diagnostic cost is included, more interaction may or may not be worth paying for.
That tradeoff is one of the main empirical targets of the paper.

\paragraph{Interaction ladder.}
The experimental family orders admissible measurements by cost and task contact:
\begin{align*}
b_0&:\text{no interaction},&
b_1&:\text{read-only task view},\\
b_2&:\text{verifier-only baseline probe},&
b_3&:\text{one-step scratch scout},\\
b_4&:\text{multi-step scratch scout with post-feedback bids}.&
\end{align*}
The ladder lets us ask whether the first nontrivial task contact already creates value, and whether model-coupled scratch interaction recovers more routing value than verifier-only measurement.

\subsection{Finite latent-mode reduction}

The main theorem is instance-level and does not require latent modes.
Modes enter only when they provide a reusable abstraction for pilots.

\begin{assumption}[Finite solver-relevant mode sufficiency]
\label{ass:finite-mode}
There exists a finite latent variable $S\in\cS$ such that for every action $a\in\cA$,
\begin{equation}
\E\big[\ell(a,X)\mid X\big]
=
\E\big[\ell(a,X)\mid S\big]
\qquad\text{almost surely.}
\end{equation}
Equivalently, there is a mode-conditioned loss matrix
\[
\ell(a,s):=\E\big[\ell(a,X)\mid S=s\big]
\]
that captures the action-level Bayes loss through the finite solver-relevant state $S$.
\end{assumption}

\begin{proposition}[Posterior decision on a mode-conditioned loss matrix]
\label{prop:mode-reduction}
Under Assumption~\ref{ass:finite-mode},
\begin{equation}
\bar L_b
=
\E\left[
\min_{a\in\cA}
\sum_{s\in\cS}
\Prb[S=s\mid W,Z_b]\,\ell(a,s)
\right].
\label{eq:mode-reduction}
\end{equation}
Therefore reusable estimates of $\ell(a,s)$ and posterior summaries of $\Prb[S\mid W,Z_b]$ suffice to score many candidate routers analytically.
\end{proposition}

\begin{proof}
By Assumption~\ref{ass:finite-mode},
\[
\E\big[\ell(a,X)\mid W,Z_b\big]
=
\sum_{s\in\cS}\Prb[S=s\mid W,Z_b] \ell(a,s),
\]
and substituting this expression into~\eqref{eq:rawlzb} yields~\eqref{eq:mode-reduction}.
\end{proof}

Proposition~\ref{prop:mode-reduction} gives a compact sufficient condition under which reusable pilot summaries can score many routers analytically.
Latent modes remain valuable when they let us reuse structured pilots, but the primary deployment object is still the action-level decision after a small admissible interaction.

\begin{assumption}[Packed certified-effort specialization]
\label{ass:packed-effort}
Suppose the finite mode $S\in\cS$ is sufficient in the sense of Assumption~\ref{ass:finite-mode}.
For each worker model $m\in\cM$ and mode $s$, let $p_m(s)>0$ be the certified per-attempt success probability under the fixed harness and let $\kappa(m)>0$ be the per-attempt deployment cost.
For a router that observes signal $Z$ and allocates probability mass $q_Z(s)$ to the mode-specific computation path, assume the certified effort surrogate in mode $s$ is proportional to
\begin{equation}
T_{m,q}(s,Z)\propto \frac{\kappa(m)}{p_m(s)\,q_Z(s)} .
\label{eq:packed-effort}
\end{equation}
\end{assumption}

This specialization is deliberately stronger than the action-level identity in Theorem~\ref{thm:main-identity}.
It is useful because it recovers the four interpretable Achille--Soatto-style terms from the older draft: cost, competence, information, and mismatch.
It should be read as an explanatory reduction, not as an assumption required for the main AutoResearch regret estimates.

\begin{definition}[Cost-adjusted certified log-effort]
\label{def:log-effort}
Under Assumption~\ref{ass:packed-effort}, define
\begin{equation}
\J_{\mathrm{pack}}(m,q)
=
\log\kappa(m)
+\E_{S\sim\pi}[-\log p_m(S)]
+\E_Z\!\left[
  \Hent(\pi_Z)+\KL(\pi_Z\|q_Z)
\right],
\label{eq:packed-objective}
\end{equation}
where $\pi_Z(s)=\Prb[S=s\mid Z]$.
Lower $\J_{\mathrm{pack}}$ means lower expected certified log-effort up to constants that do not affect design choice.
\end{definition}

\begin{theorem}[Packed-mode certified-effort decomposition]
\label{thm:packed-four-term}
Under Assumption~\ref{ass:packed-effort}, compare a prior-matched baseline design $(m_0,\pi)$ with a deployed design $(m,q)$.
The improvement in expected certified log-effort,
$\Delta=\J_{\mathrm{pack}}(m_0,\pi)-\J_{\mathrm{pack}}(m,q)$, decomposes as
\begin{equation}
\Delta
=
\underbrace{\log\frac{\kappa(m_0)}{\kappa(m)}}_{\mathrm{cost}}
+
\underbrace{\sum_{s\in\cS}\pi_s\log\frac{p_m(s)}{p_{m_0}(s)}}_{\mathrm{competence}}
+
\underbrace{\MI_\pi(S;Z)}_{\mathrm{routing\ information}}
-
\underbrace{\E_Z\!\left[\KL(\pi_Z\|q_Z)\right]}_{\mathrm{routing\ mismatch}} .
\label{eq:packed-four-term}
\end{equation}
\end{theorem}

\begin{proof}
Taking logarithms in~\eqref{eq:packed-effort} yields
$\log T_{m,q}(s,Z)=C+\log\kappa(m)-\log p_m(s)-\log q_Z(s)$.
Averaging over $S\mid Z$ turns the final term into the cross-entropy
$-\E_{S\mid Z}\log q_Z(S)=\Hent(\pi_Z)+\KL(\pi_Z\|q_Z)$.
Subtracting the prior-matched baseline cancels constants and uses
$\Hent(\pi)-\E_Z\Hent(\pi_Z)=\MI_\pi(S;Z)$.
\end{proof}

The decomposition clarifies how to interpret the operational quantities in the main campaign.
Per-step price and latency enter through $\kappa$.
Mode-conditional entry success enters through $p_m(s)$.
The verifier-backed transcript helps only through the entropy reduction $\MI(S;Z)$ that it creates before routing.
The learned router loses the remaining opportunity through mismatch or, in the instance-level analysis, through the action regret in~\eqref{eq:routed-objective}.

\begin{proposition}[Single-mode retry crossover]
\label{prop:single-mode-crossover}
On a single-mode task, let a large worker have one-shot success probability $p_l$ and cost $\kappa_l$.
Let a cheaper worker have success probability $p_s<p_l$ and cost $\kappa_s<\kappa_l$.
The number of cheap-worker retries required to match the large worker's success probability is
\begin{equation}
d_\times=\frac{\log(1-p_l)}{\log(1-p_s)} .
\end{equation}
The cheap worker is cost-favorable at this crossover whenever $d_\times\kappa_s<\kappa_l$.
\end{proposition}

This proposition is not the main experimental protocol---we route whole workers rather than only choosing retry depth---but it explains why cost-to-success, not final loss alone, is the natural objective.
A weaker worker can be a rational deployment choice if verifier-backed interaction and lower per-step cost compensate for lower one-shot competence.

\section{Reusable pilots and design selection}
\label{sec:selection}

The previous section defines the deployment objective for a fixed diagnostic interaction and a fixed router.
A practitioner, however, must choose among many designs.
We therefore define a deployment design to be a pair
\[
d=(b,r_b), \qquad b\in\cB,
\]
with objective $\J(d):=\J_b(r_b)$.
Given a finite candidate menu $\cD$, the selected design from $N$ pilot instances is denoted $\hat d_N$ and its selection regret is
\begin{equation}
\Reg_{\mathrm{sel}}(N)
:=
\E\big[\J(\hat d_N)-\min_{d\in\cD}\J(d)\big].
\label{eq:selection-regret}
\end{equation}

The experimental proposal of the paper is not to black-box race every element of $\cD$ from scratch.
Instead, the pilot logs reusable measurement transcripts and full-run counterfactual losses.
When a finite latent-mode abstraction is useful, those pilots can be compressed into mode-conditioned losses and signal posteriors; when it is not, the same records still support direct instance-level router evaluation.

\begin{algorithm}[Reusable pilot protocol for minimal-interaction routing]
\label{alg:pilot}
\normalfont\small
\begin{tabular}{@{}r p{0.86\linewidth}@{}}
\textit{Input:} & interacting-model menu $\cM$, diagnostic ladder $\cB$, candidate router family, verifier-indexed deployment loss $\ell_V$, pilot split, holdout split.\\[0.25em]
1 & Choose a finite candidate design menu $\cD$ consisting of fixed-model baselines, static-context policies, verifier-only routers, and scratch-scout routers.\\
2 & For each pilot instance and each admissible interaction $b$, execute the standardized pre-run measurement against the same verifier $V$, log $Z_b$, and record diagnostic cost $c_b$. Scratch artifacts are discarded after measurement.\\
3 & Execute the full model menu on the pilot split from the original instance state and estimate either instance-level losses directly or, under Assumption~\ref{ass:finite-mode}, the reusable mode-conditioned loss matrix $\ell(a,s)$ and posterior surrogates for $\Prb[S\mid W,Z_b]$.\\
4 & Score every candidate design $d\in\cD$ under the accounting identities of Section~\ref{sec:accounting}; keep a small finalist set and include negative controls such as shuffled transcripts.\\
5 & Evaluate finalists on holdout instances and report best fixed performance, workload opportunity, raw and net diagnostic opportunity, router regret, selection regret, and first-hit time reduction.
\end{tabular}
\end{algorithm}

The key point is amortization.
If many candidate routers share the same model/action menu and the same admissible diagnostic family, the expensive quantities are the pilot losses and probe traces, not the combinatorial number of candidate policies built on top of them.
That is why the experimental section measures selection regret under a fixed evaluation budget rather than only reporting the best observed routed policy.

\section{Experimental campaign}
\label{sec:experiments}

This section specifies the paper-facing campaign.
Each subsection is organized around three questions: what is being tested, why the test matters, and how the test is carried out.
The goal is not to hide the benchmark behind a single headline number, but to align each empirical claim with a specific quantity from Sections~\ref{sec:setup}--\ref{sec:selection}.

\subsection{Benchmark family and operational instantiation}
\label{subsec:benchmark}

The experimental family is AutoResearch-style CIFAR-10 optimization.
Each instance contains an editable starting program \texttt{train.py}, a fixed training and validation regime, a verifier based on validation loss, and a nuisance seed.
The mainline task instances are indexed as
\[
X_{f,s}=(\texttt{train.py}_{f,0},D_f,V,B_{\mathrm{full}},s),
\]
where $f$ is the task family, $s$ is the instance seed, $D_f$ is the data/regime specification, $V$ is the verifier, and $B_{\mathrm{full}}$ is the full verifier training budget.
We use three compact families that already exist in the benchmark substrate: \texttt{mlp\_flat}, \texttt{cnn\_compact}, and \texttt{resnet\_tiny}.
They intentionally vary the starting model class while keeping the artifact interface and verifier fixed.
The seed $s$ is treated as a task-instance nuisance variable and should control data subset realization, shuffle, augmentation, label noise or imbalance when used, and model initialization whenever the implementation permits a unified seed.
For the main pilot we use five seeds per family, so the full worker matrix contains $3$ families $\times$ $5$ seeds $\times$ $3$ workers $=45$ no-stop full trajectories before router evaluations.
Additional repeated LLM trajectories on a smaller subset estimate stochasticity of the agent optimization path.

The paper-facing worker menu is the Claude-family depth menu
\[
\cM=\{\texttt{haiku},\ \texttt{sonnet},\ \texttt{opus}\}.
\]
The main router is an \texttt{opus} prompt-only orchestrator controlled by the framework.
It observes a diagnostic signal, returns a JSON decision over the worker menu, and does not edit \texttt{train.py} itself.
This is scientifically cleaner than using a fully autonomous multi-agent IDE as the main router because the framework controls spawning, reset state, logging, worker prompts, and oracle counterfactuals.
Claude Code-style subagents or Codex subagents are natural ecological ablations, but the main result uses prompt-only routing to avoid confounding routing ability with hidden tool policy.

If the same protocol is repeated with OpenAI/Codex workers, the model menu can be swapped for a small realistic solver set such as
\[
\cM=\{\texttt{gpt\_5\_4\_mini},\ \texttt{gpt\_5\_3\_codex},\ \texttt{claude\_sonnet}\},
\]
with optional appendix extensions to additional variants.
The definitions do not depend on the vendor menu; only the worker costs and empirical losses change.

We evaluate a nested family of router-visible diagnostic signals.
The fixed ladder is
\[
Z_0 \subset Z_{\mathrm{file}} \subset Z_{\mathrm{base}} \subset Z_{\mathrm{scout},1} \subset Z_{\mathrm{scout},2},
\]
where $Z_0$ exposes only the task identifier and worker menu, $Z_{\mathrm{file}}$ adds the initial \texttt{train.py} and workload summary, $Z_{\mathrm{base}}$ adds a cheap unedited baseline verifier probe, and $Z_{\mathrm{scout},k}$ adds $k$ scratch optimization steps per worker followed by verifier feedback.
The baseline probe records validation loss, validation accuracy, validation bits-per-byte when applicable, training seconds, total seconds, total optimizer steps, and parse/runtime status.
The scratch-scout rungs give every candidate worker the same micro-budget on a scratch copy of \texttt{train.py}, log verifier deltas, edit validity, runtime, and pre/post bids, then reset the instance.
After observing $Z_b$, the router selects the worker that receives the full optimization budget from the original artifact state.

We also test adaptive diagnostics.
The adaptive router receives the same diagnostic action set $\cQ$ but may choose which measurement to buy under a small diagnostic budget before it selects the worker.
Allowed diagnostic actions include reading \texttt{train.py}, running a baseline probe at budget $B_{\mathrm{diag}}\in\{32,128\}$, requesting one scratch scout, requesting a second scratch scout, and stopping to route.
The adaptive policy's chosen measurements, cost, and final worker decision are logged; its net value is compared against the fixed ladder after charging diagnostic cost.
This is an experimental test of the router--verifier interface rather than a separate theoretical contribution.

\begin{figure}[t]
\placeholder{0.96\linewidth}{
Planned figure: overview of the benchmark family.
Left panel: three compact AutoResearch families with five nuisance seeds each.
Middle panel: fixed and adaptive diagnostic measurements, from no interaction to verifier-only probe to symmetric scratch scout, all producing $Z_b$ before reset.
Right panel: prompt-only Opus router chooses one worker from $\{\texttt{haiku},\texttt{sonnet},\texttt{opus}\}$ before the full expensive run and is scored by cost-to-verified-improvement.
}
\caption{Operational picture of minimal diagnostic interaction in the AutoResearch benchmark family.}
\label{fig:overview}
\end{figure}

\subsection{Operational budgets, utilities, and logging}
\label{subsec:ops}

The worker horizon is fixed at $H=20$ edit steps.
Each edit step consists of one worker reasoning/code-generation call, materialization of the proposed structured edit, and one verifier run on the resulting \texttt{train.py}.
The diagnostic verifier budget $B_{\mathrm{diag}}$ is separated from the full-run verifier training budget $B_{\mathrm{full}}$.
This separation is essential: changing the diagnostic budget should change only the information available to the router, not the verifier used to score full-run worker performance.
Before the final campaign, we run a calibration sweep over
\[
B_{\mathrm{full}}\in\{32,128,256,512\}
\]
and choose the smallest full verifier budget for which baseline measurements are stable, worker edit rankings are not dominated by noise, and runtime remains feasible.
The same sweep chooses the diagnostic budgets used by $Z_{\mathrm{base}}$ and $Z_{\mathrm{scout},k}$.
A preliminary runtime check on the existing families with $4096$ training examples and $1024$ validation examples found that at $128$ optimizer steps, \texttt{mlp\_flat}, \texttt{cnn\_compact}, and \texttt{resnet\_tiny} took approximately $0.6$, $3.0$, and $16.4$ seconds of training time respectively; the residual architecture is therefore retained as a useful slower stress family but can be lightened if the full matrix becomes infeasible.

The primary success thresholds are relative improvements
\[
\delta\in\{0.05,0.10,0.15,0.20\}.
\]
The $5\%$ threshold is the primary claim threshold; larger thresholds are reported as sensitivity because they can be too severe for short-budget AutoResearch instances.
Full trajectories always run to horizon $H$ even after the first threshold hit.
This lets us reconstruct first-hit cost while also measuring persistence and final performance.

The primary routing utility is the censored log-effort loss in~\eqref{eq:log-effort-loss}, with wall-clock cost as the main $C$ and token/API dollars as a sensitivity cost.
For each instance $X_{f,s}$, every worker $m\in\cM$ is run counterfactually under the same full-run budget and seed.
For a fixed threshold and cost convention, the oracle worker is
\[
m^\star(X_{f,s})\in\argmin_{m\in\cM}\ell_{\log,\delta,\rho}(m,X_{f,s}),
\]
and the router regret is
\[
\reg(R,b;X_{f,s})
=
\ell_{\log,\delta,\rho}(m_R,X_{f,s})
-
\min_{m\in\cM}\ell_{\log,\delta,\rho}(m,X_{f,s}).
\]
We also report final-loss regret, hit-rate regret, cost-to-success regret, and oracle-match rate because the best worker can change when the decision utility changes.

Every router decision record stores task family, instance seed, diagnostic policy, diagnostic budgets, diagnostic verifier calls, diagnostic wall time, router model, router input and output tokens, router dollar cost, router wall time, selected worker, worker ranking, confidence, success-probability estimates, expected cost-to-success estimates, and diagnostic features used.
Every worker-step record stores step index, worker model, input and output tokens, reasoning wall time, API cost, verifier training steps, verifier training seconds, verifier wall time, validation loss, validation accuracy, relative improvement, success indicators for all thresholds, parse failures, runtime failures, and whether the artifact was promoted as the visible parent.
Every run summary stores baseline loss, best loss, final loss, first-hit step for each threshold, hit indicators, total reasoning time, total verifier time, total wall-clock time, total tokens, total cost, total verifier calls, persistence after first hit, oracle-best worker under each utility, and router regret.
Uncertainty intervals are computed by nonparametric bootstrap resampling at the task-family level with seeds nested inside families; when worker-run stochasticity is explicitly replicated, the bootstrap resamples task instances first and agent-run replicates second.

The router output contract is fixed before evaluation.
For a worker-routing decision it must return exactly one selected worker in $\{\texttt{haiku},\texttt{sonnet},\texttt{opus}\}$, a complete worker ranking, per-worker success-probability estimates, per-worker expected cost-to-success estimates or null values when no success is expected, an expected final-loss ranking, a scalar confidence, a list of diagnostic features used, and a brief rationale grounded only in the allowed transcript.
The router prompt explicitly states that the router must not optimize \texttt{train.py}, must not propose edits, and must prefer the cheapest worker likely to reach the relative-improvement success criterion.
For adaptive diagnostics, the router first chooses diagnostic actions from the predeclared set $\cQ$ under the diagnostic budget and only then emits the same worker-routing JSON.

\subsection{Experiment I: does time-to-verified-success change the ranking?}
\label{subsec:exp1}

\paragraph{What is being tested.}
We test whether interacting-model rankings induced by final improvement or hit-once pass rate agree with rankings induced by time-aware deployment losses such as $\ell_{\mathrm{hit}}$ and $\ell_{\mathrm{occ}}$.
Equivalently, we ask whether benchmark racing on a scalar success metric is already sufficient, or whether deployment-aware metrics expose important heterogeneity.

\paragraph{Why this test matters.}
If rankings are stable across these objectives, then the rest of the routing story becomes less compelling.
If instead a model that looks strong under pass rate is weak under first-hit cost, wall-clock effort, token cost, or occupancy, then routing must be framed around deployment loss rather than aggregate success.
This experiment therefore justifies the objective of the paper before any router is introduced.

\paragraph{How it is tested.}
For each family--seed instance $X_{f,s}$ and every worker in $\cM$, we run full no-stop trajectories up to the common horizon $H$ and log verifier improvement, first threshold-hit time, occupancy above threshold, wall-clock, and token cost.
From these trajectories we reconstruct per-model losses under several metrics and report rankings over nuisance draws and task families.
Additional agent-run replicates on a smaller subset estimate how much ranking instability comes from LLM trajectory stochasticity rather than the task instance.
The main comparison is between final improvement, first-hit success, cost-aware first-hit loss, and occupancy-aware loss.

\begin{figure}[t]
\placeholder{0.96\linewidth}{
Planned figure: rank comparison across metrics.
Left: per-worker bar chart of first-hit success across family--seed instances.
Center: per-model bar chart of occupancy-aware loss or cost-adjusted verified effort.
Right: rank-flip summary showing where benchmark racing and deployment ranking disagree.
}
\caption{Experiment I asks whether pass rate is already enough or whether time-to-verified-success changes the ranking.}
\label{fig:metric-ranks}
\end{figure}

\subsection{Experiment II: can a minimal verifier-backed interaction reveal routing-relevant structure?}
\label{subsec:exp2}

\paragraph{What is being tested.}
We test whether a cheap diagnostic measurement produces a signal $Z_b$ that exposes model-relevant structure beyond static controls.
The central objects are the raw diagnostic oracle $\bar L_b$ and the raw diagnostic opportunity $\bar\Delta_b$.

\paragraph{Why this test matters.}
This is the heart of the minimal-interaction thesis.
If a tiny admissible interaction already reveals enough structure to change the optimal worker or reduce first-hit time, then the verifier is not only a final judge but also a diagnostic instrument.
If it does not, then routing should stay with the best fixed worker or with family-conditioned routing.

\paragraph{How it is tested.}
We compare several admissible information states drawn from the same harness: no interaction, read-only task view, verifier-only baseline probe, symmetric scratch scout, adaptive diagnostic policy, and shuffled-transcript negative controls.
The primary report is incremental: how much does each measurement add beyond information available without interacting with the verifier?
Because task family $f$ is known to the router, we report both opportunity beyond the best fixed worker and opportunity beyond a family-conditioned workload oracle.
The latter isolates value due to live verifier-backed task contact rather than static family identity.
Diagnostic-feature analyses may visualize which parts of the transcript separate regimes, but the main claim is decision-theoretic, not classification-theoretic.

\begin{figure}[t]
\placeholder{0.96\linewidth}{
Planned figure: diagnostic signal ablation.
Bars compare best fixed, family-conditioned, read-only, verifier-only, scratch-scout, adaptive-diagnostic, and shuffled-transcript policies on the same holdout split.
Primary y-axis: raw diagnostic opportunity $\bar\Delta_b$ beyond static controls.
Secondary appendix panel: optional workload-shift or mode-confusion summaries to help interpret the signal, not to define the main claim.
}
\caption{Experiment II measures how much routing-relevant structure is exposed by a cheap admissible interaction.}
\label{fig:signal-ablation}
\end{figure}

\subsection{Experiment III: how much of the oracle routing value is recoverable?}
\label{subsec:exp3}

\paragraph{What is being tested.}
We compare deployment policies on holdout instances: the best fixed worker, the family-conditioned workload oracle, the diagnostic oracle, the Opus prompt-only router, and lightweight learned routers fit from the pilot split.
This directly estimates $\Delta_W$, $\Delta_b$, and $\Reg_b(r_b)$.

\paragraph{Why this test matters.}
It separates three qualitatively different stories.
If $\Delta_W$ is large, then known family identity already explains most routing value.
If $\Delta_b\approx 0$, then even verifier-backed pre-runs do not justify routing.
If $\Delta_b>0$ and the Opus or learned-router regret is small, then a minimal interaction genuinely changes the deployment decision and reduces time-to-verified-success.
This is the main test of whether routing is worthwhile in practice.

\paragraph{How it is tested.}
We split nuisance draws into pilot and holdout partitions.
The pilot split is used to estimate action losses, calibrate the oracle workers, and fit lightweight routers from $(W,Z_b)$ to workers; the holdout split evaluates deployment loss, first-hit time, and cost-to-hit only.
The Opus prompt-only router receives the same allowed signal and returns a JSON worker choice, ranking, confidence, success-probability estimates, expected cost-to-success estimates, and a short rationale grounded in the diagnostic transcript.
The lightweight routers remain simple and interpretable---for example nearest-prototype, linear, or tree-based rules---because they ask whether the signal is usable without a frontier router.
We then report gain over best fixed worker, first-hit time reduction, oracle-match rate, and regret to the diagnostic oracle.

\begin{figure}[t]
\placeholder{0.96\linewidth}{
Planned figure: routing recovery diagram.
Plot best fixed worker, family-conditioned oracle, diagnostic oracle, Opus router, and learned router under the same deployment loss.
Annotate the bars or connecting arrows with $\Delta_b$ and $\Reg_b(r_b)$; report $\Delta_W$ to show how much value is already explained by static task-family identity.
}
\caption{Experiment III measures the fraction of oracle routing value that is actually recovered by a learned router.}
\label{fig:routing-recovery}
\end{figure}

\subsection{Experiment IV: what is the minimum useful diagnostic interaction?}
\label{subsec:exp4}

\paragraph{What is being tested.}
We test how routing value changes as the admissible interaction becomes richer and identify both $b^{\mathrm{suf}}_\eta$ and $b^{\mathrm{net}}_\star$.
In the mainline instantiation, the interaction family is the ordered measurement ladder from no interaction to read-only task views, verifier-only probes, one-step scratch scouts, and multi-step scratch scouts.
We also include the adaptive diagnostic policy from Definition~\ref{def:adaptive-diagnostic}, which can stop early or buy an additional measurement under the same diagnostic cost accounting.

\paragraph{Why this test matters.}
The paper's central operational promise is not merely that probes can help, but that one can estimate how much probe is enough.
A practitioner rarely wants the maximally informative interaction if a cheaper one already recovers most of the routing value.
This experiment turns the theory into an estimable protocol rather than a fixed hand-coded intervention.

\paragraph{How it is tested.}
Within the measurement ladder, we sweep admissible diagnostic budgets and compute both raw and net diagnostic opportunity.
The raw curve estimates how quickly routing information saturates.
The net curve discounts that value by probe cost and identifies when further interaction stops being worth paying for.
The main report highlights the smallest fixed probe that recovers a target fraction of the best raw value, the probe that maximizes net gain, and whether adaptive measurement improves net value after paying for the extra router and verifier calls it chooses.
The key comparison is whether model-coupled scratch scouts dominate verifier-only probes after their extra model-verifier information is charged for its higher cost.

\begin{figure}[t]
\placeholder{0.96\linewidth}{
Planned figure: minimum useful diagnostic interaction.
Left: raw diagnostic opportunity as a function of measurement protocol and scout budget.
Right: net diagnostic opportunity after subtracting probe cost.
Mark the minimum sufficient interaction $b^{\mathrm{suf}}_\eta$ and the economically optimal interaction $b^{\mathrm{net}}_\star$.
}
\caption{Experiment IV turns minimal diagnostic interaction into an estimable curve rather than a slogan.}
\label{fig:minimal-interaction}
\end{figure}

\subsection{Experiment V: do reusable pilots reduce selection regret?}
\label{subsec:exp5}

\paragraph{What is being tested.}
We test whether structured reusable pilots reduce selection regret relative to black-box benchmark racing under the same evaluation budget.
The quantity of interest is $\Reg_{\mathrm{sel}}(N)$ from~\eqref{eq:selection-regret}.

\paragraph{Why this test matters.}
Even if routing is valuable, the experimental protocol is unattractive if it requires racing every candidate design end to end.
The paper claims something stronger: because losses and probe signals are reusable, the builder can often explore a larger candidate menu for the same evaluation budget.
This is the bridge from theory to methodology.

\paragraph{How it is tested.}
We construct a candidate menu over fixed-model policies, family-conditioned policies, verifier-probe routers, and optional model-coupled scout routers.
Under a fixed budget of pilot instances and holdout confirmations, we compare two strategies: black-box design racing and the reusable-pilot protocol of Algorithm~\ref{alg:pilot}.
The outcome is not only the best observed design but the regret of the selected design relative to the best design in the menu.

\begin{figure}[t]
\placeholder{0.96\linewidth}{
Planned figure: selection-regret comparison.
X-axis: total evaluation budget.
Y-axis: selection regret.
Compare black-box racing against reusable-pilot selection, with separate traces for different candidate-menu sizes.
}
\caption{Experiment V evaluates whether the structured pilot protocol improves design selection efficiency.}
\label{fig:selection-regret}
\end{figure}

\section{Discussion}
\label{sec:discussion}

The three-family mainline is intentionally small.
It trades breadth for dense counterfactual measurement: every worker can be run on every family--seed instance, so the oracle worker and router regret are observable rather than inferred from a leaderboard.
The main interpretive risk is that a router could exploit static family identity rather than live verifier feedback.
We address this by reporting both the family-conditioned workload oracle and the incremental diagnostic opportunity beyond that oracle.
If a cheap diagnostic measurement still changes the best worker after conditioning on family, then the signal is genuinely diagnostic for model--verifier behavior.
If it does not, the framework still gives a useful negative result: in that deployment, model routing is not worth engineering under the admitted probe family.

\paragraph{Beyond AutoResearch.}
AutoResearch is the controlled empirical substrate, not the scope of the framework.
The same formal objects apply to other verifier-backed tasks by changing the verifier and the admissible probe family: in repository repair, $V$ is the test suite and $b$ may be a smoke-test or failing-test probe; in theorem proving, $V$ is a proof checker and $b$ may be a bounded search or tactic replay; in data-pipeline construction, $V$ is schema, unit, or validation-data compliance and $b$ may be a dry run on a small shard.
The empirical claim of this paper is established on AutoResearch because it is cheap enough to instrument densely; the broader claim is methodological: once a verifier can produce a cheap pre-run signal, the same accounting estimates whether routing on that signal is worth paying for.

The main limitation is that all value estimates are conditional on the verifier, budget, and diagnostic family.
A probe can be too weak to expose useful structure, or too expensive to be worth running even when it is informative.
The framework also assumes that pilot and holdout instances are drawn from the same deployment distribution; robustness workloads should therefore be reported as distribution-shift stress tests rather than as the primary evidence for the main claim.

\section{Conclusion}
\label{sec:conclusion}

Minimal diagnostic interaction reframes model choice for verifier-backed agents as a measurable deployment decision.
Given a verifier, a budget, and a finite menu of worker models or harness actions, the paper asks how much cheap interaction is enough to select the action for the current instance.
The answer is not a universal model ranking, but an accounting of fixed-model performance, diagnostic opportunity, probe cost, and router regret.
This makes AutoResearch a controlled testbed for a broader claim: in verifiable transductive tasks, useful routing information is the information that reduces time to a certified success after the cost of measuring it has been paid.

\bibliographystyle{plainnat}
\bibliography{references}

\end{document}
